% ECONOMICS_PROBLEM: {"id": "E42-470", "title": "Economic role and fragility of crypto, stablecoins and DeFi", "jel_code": "E42", "source_id": "OP-0097"}
% !TeX program = xelatex
\documentclass[11pt,a4paper]{article}
\usepackage[margin=23mm]{geometry}
\usepackage{fontspec}
\setmainfont{Latin Modern Roman}
\usepackage{amsmath,amssymb,mathtools}
\usepackage{xcolor,enumitem,booktabs,longtable,array,xurl,hyperref}
\definecolor{muted}{HTML}{53636C}
\hypersetup{colorlinks=true,urlcolor=blue,linkcolor=blue}
\setlength{\parindent}{0pt}
\setlength{\parskip}{5pt}
\newcommand{\R}{\mathbb R}
\newcommand{\E}{\mathbb E}
\newcommand{\N}{\mathbb N}
\newcommand{\ind}{\mathbf 1}
\newcommand{\Var}{\operatorname{Var}}
\newcommand{\Cov}{\operatorname{Cov}}
\newcommand{\supp}{\operatorname{supp}}
\DeclareMathOperator*{\argmin}{arg\,min}
\DeclareMathOperator*{\argmax}{arg\,max}
\newcommand{\entrymeta}[3]{{\sffamily\small\color{muted}\textbf{\texttt{#1}}\quad|\quad #2\par #3\par}}
\newcommand{\Problemref}[1]{\texttt{#1}}
\newcommand{\JELref}[2]{\texttt{#2} (JEL \texttt{#1})}
\begin{document}
\section*{Economic role and fragility of crypto, stablecoins and DeFi}
\textbf{Problem ID:} \texttt{E42-470}\quad \textbf{JEL:} \texttt{E42}\par
% BEGIN ECONOMICS STATEMENT
\label{problem:OP-0097}
\label{atlas:prob:F7}

\noindent\textbf{Atlas ID:} F7; \textbf{original number:} 97; \textbf{field:} Finance. \textbf{Classification:} Editorial research formulation (theoretical/empirical); not a certified unsolved theorem.

\subsection*{Definitions and assumptions}
Fix a dated protocol and institutional regime over a finite horizon $H$. Users choose feasible payments, collateralized loans and trades; validators or liquidity providers choose activity under stated incentives. Model $\theta$ includes transaction-service utility, settlement reliability, resource costs, cyber losses and an equilibrium selection. For a redeemable stablecoin, let outstanding claims be $B_t$, redemptions due be $D_t\in[0,B_t]$, liquid reserves be $L_t$ and other reserve liquidation proceeds available by the settlement deadline be $Q_t$; redemption promises require an explicit priority and enforceability rule. Algorithmic or overcollateralized tokens require their own collateral and liquidation constraints rather than this reserve identity. All welfare and loss variables are integrable. Claims and reserve values use the same promised redemption currency.

\subsection*{Formal research question}
Compare the dated decentralized or stablecoin arrangement $a^1$ with a specified feasible payment-and-credit alternative $a^0$ serving the same population. Identify
\[
 \Delta_W=\sum_h\omega_h\mathbb E_\theta[U_h(a^1)-U_h(a^0)],
 \qquad
 p_{\mathrm{fail}}=\mathbb P_\theta(L_t+Q_t<D_t
   \text{ at a required redemption date}),
\]
where $U_h$ includes service benefits, pecuniary outcomes and losses, and $\omega_h\ge0$ are stated weights. Characterize reserve, governance and collateral rules making $\Delta_W>0$ at a given failure-risk bound. Identify exposure across on-chain positions, exchanges, custodians and reserve assets without counting the same underlying collateral multiple times.

\subsection*{Interpretation and scope}
Token prices can reflect transaction-service benefits and network effects; valuing every token solely as a discounted cash-flow claim can omit the object being priced. Reserve sufficiency at one date does not guarantee liquidity under simultaneous redemptions, fire-sale discounts or inaccessible custody. A blockchain records addresses and contract states, not necessarily beneficial owners, enforceable off-chain rights or the entire balance sheet. Experimental or quasi-experimental protocol changes must be dated because users, code and regulation adapt. Specify finality and failure criteria, account for fees as transfers where appropriate, and distinguish energy or computing resource costs from token rewards. A causal efficiency claim needs a comparable counterfactual service.

\subsection*{Source audit and status}
[S1] is a protocol white paper; [S2] studies cryptocurrency arbitrage; [S3] models token adoption. Atlas link [50], [S4], identifies a 2022 overview. It does not support the atlas’s precise assertion about recent stablecoin flows and short-term government-debt markets; that assertion is left unverified. The editorial welfare-and-fragility question has no certified current open status.

\subsection*{References}

\begin{enumerate}[label={[S\arabic*]},leftmargin=*]

\item Satoshi Nakamoto (2008). \emph{Bitcoin: A Peer-to-Peer Electronic Cash System}. Technical white paper. \url{https://bitcoin.org/bitcoin.pdf}. Date/version verification: \url{https://www.metzdowd.com/pipermail/cryptography/2008-October/014810.html}.
\textit{Locator and role:} White-paper abstract and Sections 2–5; author’s original release email dated 31 October 2008; Abstract and Sections 2–5 describe signatures, timestamping, proof of work and consensus. Protocol source, not evidence that stablecoins or DeFi improve welfare. The original dated announcement verifies the release year.

\item Igor Makarov and Antoinette Schoar (2020). \emph{Trading and Arbitrage in Cryptocurrency Markets}. Journal of Financial Economics 135(2), 293–319. \url{https://www.sciencedirect.com/science/article/pii/S0304405X19301746}.
\textit{Locator and role:} Publisher or author abstract / bibliographic record; Market segmentation, price differences and cryptocurrency arbitrage. Provides background or a benchmark result; does not establish that the editorial question remains unresolved in 2026.

\item Lin William Cong, Ye Li, and Neng Wang (2021). \emph{Tokenomics: Dynamic Adoption and Valuation}. Review of Financial Studies 34(3), 1105–1155. \url{https://doi.org/10.1093/rfs/hhaa089}.
\textit{Locator and role:} Publisher or author abstract / bibliographic record; Token transaction demand, network effects and dynamic adoption; online publication 2020, issue year 2021. Provides background or a benchmark result; does not establish that the editorial question remains unresolved in 2026.

\item Igor Makarov and Antoinette Schoar (2022). \emph{Cryptocurrencies and Decentralized Finance (DeFi)}. NBER Working Paper 30006. \url{https://doi.org/10.3386/w30006}.
\textit{Locator and role:} Publisher or author abstract / bibliographic record; Overview of cryptocurrency and DeFi mechanisms and risks. The atlas’s specific assertion about recent stablecoin flows and government-debt markets is not established by the identified overview; a separate dated empirical source is required.

\end{enumerate}

\subsection*{Common conventions: Source mathematical conventions}

Unless an entry states otherwise, agent and firm sets are finite. State and action spaces are measurable, strategies and outcome maps are measurable, probability laws are specified, and expectations exist for the targets under discussion. Infinite-horizon discount factors lie in $(0,1)$ and discounted objectives are finite. Norms, tolerances and complexity input sizes must be specified for computational comparisons. Constants or primitives that appear locally are inputs, not empirical facts asserted by this catalogue.

\textbf{Identification.} A statistical model $m\in\mathcal M$ implies an observable law $P_m$ and target $\theta(m)$. At law $P$, its identified set is
\[
\Theta(P;\mathcal M)=\{\theta(m):m\in\mathcal M,\ P_m=P\}.
\]
Point identification holds when this set is a singleton, or equivalently when $P_{m_1}=P_{m_2}$ implies $\theta(m_1)=\theta(m_2)$ for all compatible models. Sharp bounds mean bounds attained, or approached when the boundary is not attained, by compatible models. Writing this set is a definition, not a solution: the substantive task is to establish informative restrictions and derive or estimate the set. An unrestricted-confounding model may give only trivial bounds.

\textbf{Inference.} Identification, estimability and valid uncertainty quantification are separate questions. Fix $\alpha\in(0,1)$. A confidence procedure $C_n$ over a declared law class $\mathcal P$ has asymptotic uniform coverage at level $1-\alpha$ when
\[
\liminf_{n\to\infty}\inf_{P\in\mathcal P}\Pr_P\{\theta(P)\in C_n\}\geq1-\alpha.
\]
For set-identified targets the coverage event must be defined separately, for example coverage of the full identified set. If an entry uses identification-indexed models instead of uniquely indexed laws, the same coverage convention is applied over those models. Pointwise limits do not establish uniform validity, and asymptotic validity does not guarantee finite-sample precision. Sample sizes, dependence structures and weak-identification sequences are part of each application.

\textbf{Counterfactuals.} $Y_i(a)$ is a potential outcome under a specified intervention, including its timing, implementation and versions. It is not $Y_i$ conditional on voluntarily choosing $a$. A full intervention vector permits interference; shorthand scalar treatment notation does not silently impose no interference. Assignment, selection, attrition, anticipation and measurement must be specified. A claim about an instrument's local effect is not automatically a population policy effect.

\textbf{Economic equilibrium.} A counterfactual model includes best responses, constraints, feasibility, market clearing, state transitions and expectations consistent with the stated information. When several equilibria exist, either a declared selector is an input or the target is set-valued. Comparisons across policy regimes must use the stated selection convention. Reduced-form relationships are not presumed invariant after an equilibrium-changing intervention.

\textbf{Optimization and welfare.} Optimization is over the stated feasible set. If attainment is not established, $\sup$ or $\inf$ and $\varepsilon$-optimal choices replace an asserted maximizer or minimizer. For example, with value $V=\sup_{a\in\mathcal A}W(a)<\infty$, an $\varepsilon$-optimal set is $\{a\in\mathcal A:W(a)\geq V-\varepsilon\}$ for $\varepsilon>0$. Benefits, costs, risk and health or environmental outcomes can be aggregated only using declared units and conversion weights. Interpersonal utility weights, the treatment of fiscal transfers and foreign profits, discounting, and the behavioral welfare criterion are explicit inputs. A comparative-static derivative additionally requires existence and appropriate differentiability of the selected counterfactual.

\textbf{Sources.} Local labels [S1], [S2], and so forth restart in every question. Locators identify the relevant abstract, model, section or result at the level actually checked; a bibliographic or abstract-level verification is not represented as inspection of an entire proof. Working papers are distinguished from later publications and changing versions. Source summaries are paraphrased. All URL checks and editorial formulations refer to the audit date stated above.
% END ECONOMICS STATEMENT
\end{document}
